\documentclass[../main.tex]{subfiles}
\begin{document}

\chapter{Compiler Background - II}
\lessoninfo{
  video={M11, lessons 1--4 (Live-Variable Analysis; SSA; Examples of Compiler Optimizations; Divergence Analysis)},
  sources={Cytron et al.\ \cite{cytron1991ssa} (SSA construction with dominance frontiers); Chaitin \cite{chaitin1982coloring} (register allocation by graph colouring); LLVM Language Reference \cite{llvm2026langref} (\texttt{phi}, \texttt{select}); LLVM passes \cite{llvm2026passes} (\texttt{mem2reg}); PTX ISA \cite{nvidia2026ptx}},
  textbook={Aho et al.\ \cite{aho2006compilers} §6.2 Three-Address Code, §8.8 Register Allocation and Assignment, §9.1 The Principal Sources of Optimization, §9.2 Introduction to Data-Flow Analysis, §9.6 Loops in Flow Graphs; CPG \cite{nvidia2024cudaguide} §5.2.3 Multiprocessor Level, §5.3.2 Device Memory Accesses, §5.4.2 Control Flow Instructions, §7.1.5 \_\_noinline\_\_ and \_\_forceinline\_\_, §7.39 \#pragma unroll; PMPP \cite{hwu2022pmpp} ch.~10},
  keywords={live variable, liveness analysis, backward analysis, def and use sets, register allocation, interference graph, register spill, static single assignment, phi function, path-convergence criterion, dominance frontier, loop unrolling, function inlining, dead code elimination, constant propagation, strength reduction, loop-invariant code motion, induction variable, divergence analysis, thread-dependent and uniform values},
}

Lesson~10 followed a CUDA program through the compiler to PTX, cut the
PTX into basic blocks and a control-flow graph, and set up data-flow
analysis with reaching definitions, a forward problem.  This lesson
completes the toolkit with its backward counterpart, live-variable
analysis, and with its most important client, register allocation,
which on a GPU decides how many warps fit on an SM.  It then presents
static single assignment form, the representation in which LLVM, and
with it NVIDIA's LLVM-based NVVM compiler, carries a program, and the classic optimisations
that the analyses enable: loop unrolling, inlining, dead code
elimination, constant propagation, strength reduction and code motion.
The last section turns to a question that only parallel code poses:
which branches can diverge.  It connects the hardware of Lesson~6, the
SIMT stack and its reconvergence points, to an analysis the compiler
performs on the def-use chains of the thread index.

\section{Live-variable analysis}
\label{sec:l11-liveness}

Reaching definitions asks, for each program point, which assignments
may have produced the values seen there.\source{M11, liveness; Aho et
al.\ §9.2 Introduction to Data-Flow Analysis \cite{aho2006compilers}.}
Liveness asks the opposite question: which of the values present at a
point will still be needed later.

\begin{definition}[Live variable]
  \label{def:l11-live}
  A variable $x$ is a \term{live variable}[生存変数] at a program point
  $p$ if the value of $x$ at $p$ may be used along some path in the
  control-flow graph that starts at $p$, before $x$ is redefined on that
  path; otherwise $x$ is \emph{dead} at $p$.  The data-flow analysis
  that computes, for every point, the set of live variables is
  \term{liveness analysis}[生存変数解析].
\end{definition}

\margin{Lesson~6 used the same notion for the hardware: a live register
holds a value that will still be read.  Liveness analysis computes it
at compile time for every variable and every point.}%
A dead value occupies a register for nothing.  The main use of liveness
is therefore register allocation (\cref{sec:l11-regalloc}): only live
variables need registers, and the register of a variable that has died
can be handed to another.  The same information tells which
assignments are useless, the basis of the dead code elimination of
\cref{sec:l11-optimizations}.

\subsection{Data-flow equations}

The effect of one basic block $B$ on liveness is summarised by two sets.

\begin{definition}[Def and use sets]
  \label{def:l11-def-use}
  The \term{def set}[def集合] $\mathit{def}_B$ of a block $B$ is the set
  of variables that $B$ defines before any use of them in $B$.  The
  \term{use set}[use集合] $\mathit{use}_B$ is the set of variables that
  $B$ may use before any definition of them in $B$.
\end{definition}

\caution{An assignment such as \texttt{k = k - 1} reads \texttt{k}
before writing it, so \texttt{k} belongs to $\mathit{use}_B$ and, by
the definition, not to $\mathit{def}_B$.  Listing it in both sets, as
is sometimes done, gives the same IN set, because a variable in
$\mathit{use}_B$ is live on entry anyway.}%
A variable in $\mathit{use}_B$ is live on entry to $B$ whatever happens
after $B$, since $B$ itself reads it.  A variable in $\mathit{def}_B$ is
dead on entry to $B$ unless $B$ reads it first, since $B$ overwrites it.
Every other variable passes through $B$ unchanged: it is live on entry
exactly when it is live on exit.  With IN$[B]$ and OUT$[B]$ the sets of
live variables just before and just after $B$, this gives the equations
\begin{align}
  \mathrm{IN}[\mathrm{EXIT}] &= \varnothing, \label{eq:l11-boundary}\\
  \mathrm{IN}[B] &= \mathit{use}_B \cup
     \bigl(\mathrm{OUT}[B] - \mathit{def}_B\bigr), \label{eq:l11-in}\\
  \mathrm{OUT}[B] &= \bigcup_{S \in \mathrm{succ}(B)} \mathrm{IN}[S].
     \label{eq:l11-out}
\end{align}
\margin{Notation as in Lesson~10: $\mathrm{succ}(B)$ is the set of
successors of $B$, IN and OUT the values at the program points before
and after the block.}%
\Cref{eq:l11-boundary} is the boundary condition: no variable is live
when the program ends.  \Cref{eq:l11-in} is the transfer function of
the block, \cref{eq:l11-out} the meet: a variable is live on exit from
$B$ if and only if it is live on entry to at least one successor, since
some path continues through that successor.

\begin{definition}[Backward analysis]
  \label{def:l11-backward}
  A data-flow problem in which information flows against the edges of
  the control-flow graph, the IN set of a block being computed from its
  OUT set and the OUT set from the IN sets of the successors, is a
  \term{backward analysis}[後ろ向き解析].  Liveness is backward because
  whether a value is needed depends on what the program does
  \emph{later}; reaching definitions is forward because whether a
  definition arrives depends on what happened \emph{earlier}.
\end{definition}

The equations are solved by the same iterative method as reaching
definitions, run in the other direction.  Every IN set starts empty
(\cref{eq:l11-boundary} and the initial guess for every other block);
then the blocks are visited repeatedly, each time computing OUT$[B]$
from the current IN sets of the successors and IN$[B]$ from OUT$[B]$,
until a whole pass changes no IN set.\tip{Visit the blocks in reverse
order (exit first) for a backward problem: most values are then
already up to date when they are read, and fewer passes are needed.
The result does not depend on the order, only the work does.}  The
sets only grow, and there are finitely many variables, so the method
reaches a fixed point; it is the smallest solution of the equations,
which is the precise one.

\begin{figure}[tb]
  \centering
  % Live-variable analysis (Lesson 11): the control-flow graph of the
% chapter's own running example, a loop that accumulates s while k
% counts down, with an inner branch on t.  Each block shows its code
% and its use and def sets; the back edge B4 -> B2 closes the loop.
% Original figure (replaces the lecture's m11l1_live_eg, whose CFG the
% book does not reuse).
% Usage: % Live-variable analysis (Lesson 11): the control-flow graph of the
% chapter's own running example, a loop that accumulates s while k
% counts down, with an inner branch on t.  Each block shows its code
% and its use and def sets; the back edge B4 -> B2 closes the loop.
% Original figure (replaces the lecture's m11l1_live_eg, whose CFG the
% book does not reuse).
% Usage: % Live-variable analysis (Lesson 11): the control-flow graph of the
% chapter's own running example, a loop that accumulates s while k
% counts down, with an inner branch on t.  Each block shows its code
% and its use and def sets; the back edge B4 -> B2 closes the loop.
% Original figure (replaces the lecture's m11l1_live_eg, whose CFG the
% book does not reuse).
% Usage: \input{figures/l11-liveness-cfg}
% Width: body (120 mm); the drawing is about 104 mm wide.
\begin{tikzpicture}[x=1mm, y=1mm, font=\scriptsize\sffamily,
    >={Stealth[length=1.6mm]},
    bb/.style={draw, fill=ln-box, minimum width=40mm, align=left,
               inner sep=2pt, rounded corners=0.8mm},
    term/.style={draw, rounded corners=2mm, minimum width=14mm,
                 minimum height=5mm, inner sep=1pt}]
  \node[term] (en) at (36,100) {ENTRY};
  \node[bb] (b1) at (36,89)
    {\textbf{B1}\quad\texttt{s = 0}\\
     \hphantom{\textbf{B1}\quad}\texttt{k = n}\\[1pt]
     {\color{ln-accent}use $\{n\}$\quad def $\{s,k\}$}};
  \node[bb] (b2) at (36,71)
    {\textbf{B2}\quad\texttt{k = k - 1}\\
     \hphantom{\textbf{B2}\quad}\texttt{t = k * w}\\
     \hphantom{\textbf{B2}\quad}\texttt{if (t > 0)}\\[1pt]
     {\color{ln-accent}use $\{k,w\}$\quad def $\{t\}$}};
  \node[bb] (b3) at (86,58)
    {\textbf{B3}\quad\texttt{s = s + t}\\[1pt]
     {\color{ln-accent}use $\{s,t\}$\quad def $\varnothing$}};
  \node[bb] (b4) at (36,47)
    {\textbf{B4}\quad\texttt{if (k > 0)}\\[1pt]
     {\color{ln-accent}use $\{k\}$\quad def $\varnothing$}};
  \node[bb] (b5) at (36,33)
    {\textbf{B5}\quad\texttt{return s + v}\\[1pt]
     {\color{ln-accent}use $\{s,v\}$\quad def $\varnothing$}};
  \node[term] (ex) at (36,22) {EXIT};
  \draw[->] (en) -- (b1);
  \draw[->] (b1) -- (b2);
  \draw[->] (b2.east) -| node[pos=0.25, above] {\texttt{t > 0}} (b3.north);
  \draw[->] (b2) -- node[right] {\texttt{t <= 0}} (b4);
  \draw[->] (b3.south) |- (b4.east);
  \draw[->] (b4) -- node[right] {\texttt{k <= 0}} (b5);
  \draw[->] (b5) -- (ex);
  \draw[->, ln-caution, thick] (b4.west) -- ++(-6,0) |- 
    node[pos=0.25, left, align=right] {\iflangja{後退辺}{back\\edge}\\\texttt{k > 0}} (b2.west);
  \node[anchor=west, text=ln-muted, align=left] at (68,89)
    {\iflangja{入力: $n, w, v$}{inputs: $n, w, v$}};
\end{tikzpicture}

% Width: body (120 mm); the drawing is about 104 mm wide.
\begin{tikzpicture}[x=1mm, y=1mm, font=\scriptsize\sffamily,
    >={Stealth[length=1.6mm]},
    bb/.style={draw, fill=ln-box, minimum width=40mm, align=left,
               inner sep=2pt, rounded corners=0.8mm},
    term/.style={draw, rounded corners=2mm, minimum width=14mm,
                 minimum height=5mm, inner sep=1pt}]
  \node[term] (en) at (36,100) {ENTRY};
  \node[bb] (b1) at (36,89)
    {\textbf{B1}\quad\texttt{s = 0}\\
     \hphantom{\textbf{B1}\quad}\texttt{k = n}\\[1pt]
     {\color{ln-accent}use $\{n\}$\quad def $\{s,k\}$}};
  \node[bb] (b2) at (36,71)
    {\textbf{B2}\quad\texttt{k = k - 1}\\
     \hphantom{\textbf{B2}\quad}\texttt{t = k * w}\\
     \hphantom{\textbf{B2}\quad}\texttt{if (t > 0)}\\[1pt]
     {\color{ln-accent}use $\{k,w\}$\quad def $\{t\}$}};
  \node[bb] (b3) at (86,58)
    {\textbf{B3}\quad\texttt{s = s + t}\\[1pt]
     {\color{ln-accent}use $\{s,t\}$\quad def $\varnothing$}};
  \node[bb] (b4) at (36,47)
    {\textbf{B4}\quad\texttt{if (k > 0)}\\[1pt]
     {\color{ln-accent}use $\{k\}$\quad def $\varnothing$}};
  \node[bb] (b5) at (36,33)
    {\textbf{B5}\quad\texttt{return s + v}\\[1pt]
     {\color{ln-accent}use $\{s,v\}$\quad def $\varnothing$}};
  \node[term] (ex) at (36,22) {EXIT};
  \draw[->] (en) -- (b1);
  \draw[->] (b1) -- (b2);
  \draw[->] (b2.east) -| node[pos=0.25, above] {\texttt{t > 0}} (b3.north);
  \draw[->] (b2) -- node[right] {\texttt{t <= 0}} (b4);
  \draw[->] (b3.south) |- (b4.east);
  \draw[->] (b4) -- node[right] {\texttt{k <= 0}} (b5);
  \draw[->] (b5) -- (ex);
  \draw[->, ln-caution, thick] (b4.west) -- ++(-6,0) |- 
    node[pos=0.25, left, align=right] {\iflangja{後退辺}{back\\edge}\\\texttt{k > 0}} (b2.west);
  \node[anchor=west, text=ln-muted, align=left] at (68,89)
    {\iflangja{入力: $n, w, v$}{inputs: $n, w, v$}};
\end{tikzpicture}

% Width: body (120 mm); the drawing is about 104 mm wide.
\begin{tikzpicture}[x=1mm, y=1mm, font=\scriptsize\sffamily,
    >={Stealth[length=1.6mm]},
    bb/.style={draw, fill=ln-box, minimum width=40mm, align=left,
               inner sep=2pt, rounded corners=0.8mm},
    term/.style={draw, rounded corners=2mm, minimum width=14mm,
                 minimum height=5mm, inner sep=1pt}]
  \node[term] (en) at (36,100) {ENTRY};
  \node[bb] (b1) at (36,89)
    {\textbf{B1}\quad\texttt{s = 0}\\
     \hphantom{\textbf{B1}\quad}\texttt{k = n}\\[1pt]
     {\color{ln-accent}use $\{n\}$\quad def $\{s,k\}$}};
  \node[bb] (b2) at (36,71)
    {\textbf{B2}\quad\texttt{k = k - 1}\\
     \hphantom{\textbf{B2}\quad}\texttt{t = k * w}\\
     \hphantom{\textbf{B2}\quad}\texttt{if (t > 0)}\\[1pt]
     {\color{ln-accent}use $\{k,w\}$\quad def $\{t\}$}};
  \node[bb] (b3) at (86,58)
    {\textbf{B3}\quad\texttt{s = s + t}\\[1pt]
     {\color{ln-accent}use $\{s,t\}$\quad def $\varnothing$}};
  \node[bb] (b4) at (36,47)
    {\textbf{B4}\quad\texttt{if (k > 0)}\\[1pt]
     {\color{ln-accent}use $\{k\}$\quad def $\varnothing$}};
  \node[bb] (b5) at (36,33)
    {\textbf{B5}\quad\texttt{return s + v}\\[1pt]
     {\color{ln-accent}use $\{s,v\}$\quad def $\varnothing$}};
  \node[term] (ex) at (36,22) {EXIT};
  \draw[->] (en) -- (b1);
  \draw[->] (b1) -- (b2);
  \draw[->] (b2.east) -| node[pos=0.25, above] {\texttt{t > 0}} (b3.north);
  \draw[->] (b2) -- node[right] {\texttt{t <= 0}} (b4);
  \draw[->] (b3.south) |- (b4.east);
  \draw[->] (b4) -- node[right] {\texttt{k <= 0}} (b5);
  \draw[->] (b5) -- (ex);
  \draw[->, ln-caution, thick] (b4.west) -- ++(-6,0) |- 
    node[pos=0.25, left, align=right] {\iflangja{後退辺}{back\\edge}\\\texttt{k > 0}} (b2.west);
  \node[anchor=west, text=ln-muted, align=left] at (68,89)
    {\iflangja{入力: $n, w, v$}{inputs: $n, w, v$}};
\end{tikzpicture}

  \caption{The running example of this section: a loop over $k$ with a
    branch on $t$.  Each block lists its code and its use and def sets;
    $n$, $w$ and $v$ are inputs.}
  \label{fig:l11-liveness-cfg}
\end{figure}

\begin{example}
  \label{ex:l11-liveness}
  \Cref{fig:l11-liveness-cfg} shows a loop that counts $k$ down from
  $n$ and adds $t = k\cdot w$ to $s$ whenever $t$ is positive; the
  result is $s + v$.  The blocks are visited in the order B5, B4, B3,
  B2, B1, and \cref{tab:l11-liveness} gives the sets after the first
  and second passes.
  \begin{itemize}
    \item \emph{Pass 1.}  OUT[B5] $= \mathrm{IN}[\mathrm{EXIT}] =
      \varnothing$, so IN[B5] $= \{s,v\}$.  B4 has the successors B2
      and B5; IN[B2] is still empty, so OUT[B4] $= \{s,v\}$ and
      IN[B4] $= \{k\} \cup \{s,v\} = \{k,s,v\}$.  Then IN[B3] $=
      \{s,t\} \cup \{k,s,v\} = \{k,s,t,v\}$.  OUT[B2] is the union of
      IN[B3] and IN[B4], $\{k,s,t,v\}$; removing $\mathit{def} =
      \{t\}$ and adding $\mathit{use} = \{k,w\}$ gives IN[B2] $=
      \{k,s,v,w\}$.  Finally IN[B1] $= \{n\} \cup (\{k,s,v,w\} -
      \{s,k\}) = \{n,v,w\}$.
    \item \emph{Pass 2.}  The back edge now carries information: IN[B2]
      is no longer empty, so OUT[B4] $= \{k,s,v,w\}$ and $w$ enters
      IN[B4] and IN[B3].  OUT[B2] gains $w$, but $w$ was already in
      IN[B2], so IN[B2] and everything above it are unchanged.
    \item \emph{Pass 3} recomputes every set and changes nothing: the
      fixed point is reached.
  \end{itemize}
  The result is what the code suggests: the inputs $w$ and $v$ are live
  throughout the loop because they are read in every iteration and at
  the end; $n$ is live only on entry to B1; and $t$ is live only on the
  edge from B2 to B3, so its register is free again in B4.\margin{Our
  numbers: five variables of the loop body live at once at most
  ($k, s, t, v, w$ on the edge B2$\to$B3), so the loop needs five
  registers before any spilling; $n$ can share one of them, being dead
  once B1 has copied it into $k$.}
\end{example}

\begin{table}[tb]
  \centering\small
  \setlength{\tabcolsep}{4pt}
  \begin{tabular}{@{}lllllll@{}}
    \toprule
    Block & $\mathit{use}_B$ & $\mathit{def}_B$ & OUT$^1$ & IN$^1$ & OUT$^2$ & IN$^2$ \\
    \midrule
    B5 & $\{s,v\}$ & $\varnothing$ & $\varnothing$ & $\{s,v\}$ & $\varnothing$ & $\{s,v\}$ \\
    B4 & $\{k\}$ & $\varnothing$ & $\{s,v\}$ & $\{k,s,v\}$ & $\{k,s,v,w\}$ & $\{k,s,v,w\}$ \\
    B3 & $\{s,t\}$ & $\varnothing$ & $\{k,s,v\}$ & $\{k,s,t,v\}$ & $\{k,s,v,w\}$ & $\{k,s,t,v,w\}$ \\
    B2 & $\{k,w\}$ & $\{t\}$ & $\{k,s,t,v\}$ & $\{k,s,v,w\}$ & $\{k,s,t,v,w\}$ & $\{k,s,v,w\}$ \\
    B1 & $\{n\}$ & $\{s,k\}$ & $\{k,s,v,w\}$ & $\{n,v,w\}$ & $\{k,s,v,w\}$ & $\{n,v,w\}$ \\
    \bottomrule
  \end{tabular}
  \caption{Liveness for \cref{fig:l11-liveness-cfg}, blocks in the
    order visited.  Superscripts number the passes; the third pass
    repeats the second.}
  \label{tab:l11-liveness}
\end{table}

\subsection{Register allocation}
\label{sec:l11-regalloc}

\begin{definition}[Register allocation]
  \label{def:l11-regalloc}
  The compiler pass that maps the values of a program onto the finite
  set of machine registers is \term{register allocation}[レジスタ割り付け].
  Two values may share a register if they are never live at the same
  point.  The \term{interference graph}[干渉グラフ] has one node per
  value and an edge between two values that are live at the same time;
  an assignment of registers is a colouring of this graph in which
  adjacent nodes get different colours.
\end{definition}

\margin{Graph-colouring register allocation, with spilling when the
colouring fails, goes back to Chaitin \cite{chaitin1982coloring}; Aho
et al.\ §8.8 \cite{aho2006compilers} presents it as the standard
global method.}%
\Cref{fig:l11-interference} applies this to a straight-line sequence.
After instruction~2 neither $p$ nor $a$ is read again, so the register
of either can receive $b$; after instruction~4 neither $b$ nor $c$ is
read again, so the register of either can receive $d$.  Only two pairs
of values are ever live together, $p$ with $a$ and $b$ with $c$, so
two registers serve five values.

The same reasoning on a branching program uses the IN and OUT sets of
the blocks.  In \cref{fig:l11-liveness-cfg}, $t$ is not in OUT[B4]: of
the two successors of B4, B2 redefines $t$ before reading it and B5
never reads it, so $t$ is dead on the back edge; its register is free
for other values from the last read of $t$ until B2 writes it again.
Conversely, B2 and B4 never read $s$, yet $s$ is in their IN and OUT
sets (\cref{tab:l11-liveness}): B3 and B5 read it later, so $s$ stays
live on every path to those reads and keeps its register throughout
the loop.

\begin{figure}[tb]
  \centering
  % Register allocation from liveness (Lesson 11).  Left: a five-
% instruction sequence and the live range of each of its five values
% (bar from the definition to the last use).  Right: the interference
% graph (an edge joins two values that are live at the same time) and
% a colouring with two registers, R1 and R2.
% Original figure.
% Usage: % Register allocation from liveness (Lesson 11).  Left: a five-
% instruction sequence and the live range of each of its five values
% (bar from the definition to the last use).  Right: the interference
% graph (an edge joins two values that are live at the same time) and
% a colouring with two registers, R1 and R2.
% Original figure.
% Usage: % Register allocation from liveness (Lesson 11).  Left: a five-
% instruction sequence and the live range of each of its five values
% (bar from the definition to the last use).  Right: the interference
% graph (an edge joins two values that are live at the same time) and
% a colouring with two registers, R1 and R2.
% Original figure.
% Usage: \input{figures/l11-interference}
% Width: body (120 mm); the drawing is about 112 mm wide.
\begin{tikzpicture}[x=1mm, y=1mm, font=\scriptsize\sffamily,
    vr/.style={circle, draw, minimum size=7mm, inner sep=0pt}]
  % instructions
  \foreach \i/\code in {1/{a = p + 1}, 2/{b = a * p}, 3/{c = b + 1},
                        4/{d = c * b}, 5/{return d}} {
    \node[anchor=west] at (0,46-\i*8) {\i};
    \node[anchor=west, font=\scriptsize\ttfamily] at (4,46-\i*8) {\code};
  }
  % variable columns
  \foreach \v/\x in {p/34, a/41, b/48, c/55, d/62}
    \node[font=\scriptsize\ttfamily\bfseries] at (\x,44) {\v};
  % live ranges: from just below the defining row (or the top for the
  % input p) to the row of the last use
  \foreach \x/\ya/\yb/\col in {34/41/30/ln-accent, 41/37/30/ln-tip,
                               48/29/14/ln-accent, 55/21/14/ln-tip,
                               62/13/6/ln-accent} {
    \fill[\col!60] (\x-1.2,\yb) rectangle (\x+1.2,\ya);
  }
  \draw[ln-rule] (-1,41.5) -- (66,41.5);
  \node[anchor=north, text=ln-muted, align=center] at (48,3)
    {\iflangja{棒 = 生存範囲（定義から最後の使用まで）}{bar = live range (definition to last use)}};
  % interference graph
  \node[anchor=west, font=\scriptsize\sffamily\bfseries] at (74,44)
    {\iflangja{干渉グラフ}{interference graph}};
  \node[vr, fill=ln-accent!30] (p) at (80,32) {\texttt{p}};
  \node[vr, fill=ln-tip!30]    (a) at (96,32) {\texttt{a}};
  \node[vr, fill=ln-accent!30] (b) at (80,18) {\texttt{b}};
  \node[vr, fill=ln-tip!30]    (c) at (96,18) {\texttt{c}};
  \node[vr, fill=ln-accent!30] (d) at (88,6)  {\texttt{d}};
  \draw[thick] (p) -- (a);
  \draw[thick] (b) -- (c);
  \node[anchor=west, align=left] at (104,32)
    {\colorbox{ln-accent!30}{R1}\\\texttt{p, b, d}};
  \node[anchor=west, align=left] at (104,18)
    {\colorbox{ln-tip!30}{R2}\\\texttt{a, c}};
\end{tikzpicture}

% Width: body (120 mm); the drawing is about 112 mm wide.
\begin{tikzpicture}[x=1mm, y=1mm, font=\scriptsize\sffamily,
    vr/.style={circle, draw, minimum size=7mm, inner sep=0pt}]
  % instructions
  \foreach \i/\code in {1/{a = p + 1}, 2/{b = a * p}, 3/{c = b + 1},
                        4/{d = c * b}, 5/{return d}} {
    \node[anchor=west] at (0,46-\i*8) {\i};
    \node[anchor=west, font=\scriptsize\ttfamily] at (4,46-\i*8) {\code};
  }
  % variable columns
  \foreach \v/\x in {p/34, a/41, b/48, c/55, d/62}
    \node[font=\scriptsize\ttfamily\bfseries] at (\x,44) {\v};
  % live ranges: from just below the defining row (or the top for the
  % input p) to the row of the last use
  \foreach \x/\ya/\yb/\col in {34/41/30/ln-accent, 41/37/30/ln-tip,
                               48/29/14/ln-accent, 55/21/14/ln-tip,
                               62/13/6/ln-accent} {
    \fill[\col!60] (\x-1.2,\yb) rectangle (\x+1.2,\ya);
  }
  \draw[ln-rule] (-1,41.5) -- (66,41.5);
  \node[anchor=north, text=ln-muted, align=center] at (48,3)
    {\iflangja{棒 = 生存範囲（定義から最後の使用まで）}{bar = live range (definition to last use)}};
  % interference graph
  \node[anchor=west, font=\scriptsize\sffamily\bfseries] at (74,44)
    {\iflangja{干渉グラフ}{interference graph}};
  \node[vr, fill=ln-accent!30] (p) at (80,32) {\texttt{p}};
  \node[vr, fill=ln-tip!30]    (a) at (96,32) {\texttt{a}};
  \node[vr, fill=ln-accent!30] (b) at (80,18) {\texttt{b}};
  \node[vr, fill=ln-tip!30]    (c) at (96,18) {\texttt{c}};
  \node[vr, fill=ln-accent!30] (d) at (88,6)  {\texttt{d}};
  \draw[thick] (p) -- (a);
  \draw[thick] (b) -- (c);
  \node[anchor=west, align=left] at (104,32)
    {\colorbox{ln-accent!30}{R1}\\\texttt{p, b, d}};
  \node[anchor=west, align=left] at (104,18)
    {\colorbox{ln-tip!30}{R2}\\\texttt{a, c}};
\end{tikzpicture}

% Width: body (120 mm); the drawing is about 112 mm wide.
\begin{tikzpicture}[x=1mm, y=1mm, font=\scriptsize\sffamily,
    vr/.style={circle, draw, minimum size=7mm, inner sep=0pt}]
  % instructions
  \foreach \i/\code in {1/{a = p + 1}, 2/{b = a * p}, 3/{c = b + 1},
                        4/{d = c * b}, 5/{return d}} {
    \node[anchor=west] at (0,46-\i*8) {\i};
    \node[anchor=west, font=\scriptsize\ttfamily] at (4,46-\i*8) {\code};
  }
  % variable columns
  \foreach \v/\x in {p/34, a/41, b/48, c/55, d/62}
    \node[font=\scriptsize\ttfamily\bfseries] at (\x,44) {\v};
  % live ranges: from just below the defining row (or the top for the
  % input p) to the row of the last use
  \foreach \x/\ya/\yb/\col in {34/41/30/ln-accent, 41/37/30/ln-tip,
                               48/29/14/ln-accent, 55/21/14/ln-tip,
                               62/13/6/ln-accent} {
    \fill[\col!60] (\x-1.2,\yb) rectangle (\x+1.2,\ya);
  }
  \draw[ln-rule] (-1,41.5) -- (66,41.5);
  \node[anchor=north, text=ln-muted, align=center] at (48,3)
    {\iflangja{棒 = 生存範囲（定義から最後の使用まで）}{bar = live range (definition to last use)}};
  % interference graph
  \node[anchor=west, font=\scriptsize\sffamily\bfseries] at (74,44)
    {\iflangja{干渉グラフ}{interference graph}};
  \node[vr, fill=ln-accent!30] (p) at (80,32) {\texttt{p}};
  \node[vr, fill=ln-tip!30]    (a) at (96,32) {\texttt{a}};
  \node[vr, fill=ln-accent!30] (b) at (80,18) {\texttt{b}};
  \node[vr, fill=ln-tip!30]    (c) at (96,18) {\texttt{c}};
  \node[vr, fill=ln-accent!30] (d) at (88,6)  {\texttt{d}};
  \draw[thick] (p) -- (a);
  \draw[thick] (b) -- (c);
  \node[anchor=west, align=left] at (104,32)
    {\colorbox{ln-accent!30}{R1}\\\texttt{p, b, d}};
  \node[anchor=west, align=left] at (104,18)
    {\colorbox{ln-tip!30}{R2}\\\texttt{a, c}};
\end{tikzpicture}

  \caption{Liveness drives register allocation.  Left: the live range
    of each value of a five-instruction sequence.  Right: the
    interference graph; two colours suffice, so the five values fit
    into two registers.}
  \label{fig:l11-interference}
\end{figure}

When more values are live at some point than there are registers, the
colouring fails and some values must leave the register file.

\begin{definition}[Spill and fill]
  \label{def:l11-spill}
  A \term{register spill}[レジスタスピル] stores a live value from a
  register to memory, usually a stack slot, to free the register; the
  later load that brings the value back into a register before its next
  use is a \term{fill (reload)}[フィル（再ロード）].
\end{definition}

PTX has no register allocation: its registers are virtual and
unlimited in number, and the PTX that the front end emits contains no
spills.\margin{A PTX kernel declares as many virtual registers as it
needs, e.g.\ \texttt{.reg .b32 \%r<100>;} for \texttt{\%r0} to
\texttt{\%r99} (PTX ISA, Parameterized Variable Names
\cite{nvidia2026ptx}).}  Registers are allocated by \texttt{ptxas} when
it translates PTX into SASS for a real architecture, and the spills it
inserts go to the thread's local memory (Lesson~3).  CPG §5.3.2 lists
\enquote{any variable if the kernel uses more registers than available
(this is also known as register spilling)} among the variables placed
in local memory \cite{nvidia2024cudaguide}.  Local memory is private
to a thread but resides in device memory, so a fill that misses in the
L1 cache costs as much as a global memory access.

On a GPU the allocator balances two costs.  Fewer registers per thread
mean more resident warps, and so more latency hiding (Lesson~3); more
registers mean fewer spills.  A single register can matter: on an SM
with 65,536 registers, two blocks of 512 threads fit at 64 registers per
thread, but only one at 65 (CPG §5.2.3 Multiprocessor Level
\cite{nvidia2024cudaguide}).
The limit is 255 registers per thread (CPG Table~21
\cite{nvidia2024cudaguide}), and the programmer can
lower it with the \texttt{-maxrregcount} option or the
\texttt{\_\_launch\_bounds\_\_()} qualifier, accepting spills for
occupancy.\tip{Compile with \mbox{\texttt{--ptxas-options=-v}}: the compiler then
reports, per kernel, its registers and its local memory use
(\texttt{lmem}, CPG §5.3.2 \cite{nvidia2024cudaguide}); \texttt{ptxas}
also prints the bytes of spill stores and loads, and nonzero values
mean the allocator ran out of registers.}

\subsection{Liveness and reaching definitions compared}

The two analyses are the standard pair of data-flow problems: the same
iterative solver, different sets and opposite directions
(\cref{tab:l11-compare}).

\begin{table}[tb]
  \centering\small
  \begin{tabular}{@{}>{\raggedright\arraybackslash}p{24mm}
                    >{\raggedright\arraybackslash}p{42mm}
                    >{\raggedright\arraybackslash}p{42mm}@{}}
    \toprule
    & Reaching definitions (Lesson~10) & Live variables \\
    \midrule
    Domain & sets of definitions & sets of variables \\
    Direction & forward & backward \\
    Transfer function & $\mathit{gen}_B \cup (x - \mathit{kill}_B)$ &
      $\mathit{use}_B \cup (x - \mathit{def}_B)$ \\
    Meet & $\cup$ over predecessors & $\cup$ over successors \\
    Equations & OUT$[B] = f_B(\mathrm{IN}[B])$, IN$[B] = \bigcup_{P}
      \mathrm{OUT}[P]$ & IN$[B] = f_B(\mathrm{OUT}[B])$, OUT$[B] =
      \bigcup_{S} \mathrm{IN}[S]$ \\
    Boundary & OUT[ENTRY] $= \varnothing$ & IN[EXIT] $= \varnothing$ \\
    Initialisation & OUT$[B] = \varnothing$ & IN$[B] = \varnothing$ \\
    \bottomrule
  \end{tabular}
  \caption{Reaching definitions and liveness side by side.}
  \label{tab:l11-compare}
\end{table}

\begin{keypoint}
  Liveness runs backward from IN[EXIT] $= \varnothing$ with
  IN$[B] = \mathit{use}_B \cup (\mathrm{OUT}[B] - \mathit{def}_B)$ and
  OUT$[B]$ the union over the successors.  Values that are never live
  together can share a register; when too many are live at once,
  \texttt{ptxas} spills some to local memory.
\end{keypoint}

\section{Static single assignment}
\label{sec:l11-ssa}

In ordinary three-address code a variable may be assigned many times,
so every analysis must first find out which assignment a use
refers to, which is the question reaching definitions
answers.\source{M11, SSA; Aho et al.\ §6.2 Three-Address Code
\cite{aho2006compilers}; Cytron et al.\ \cite{cytron1991ssa}.}  The link
from a definition to the uses it reaches is a
\term{def-use chain}[定義・使用連鎖].  Static single assignment form
builds the chains into the names of the program.

\begin{definition}[Static single assignment]
  \label{def:l11-ssa}
  A program is in \term{SSA form}[SSA形式] when every variable is
  assigned at exactly one place in the program text.  The property is
  called \term{static single assignment (SSA)}[静的単一代入]:
  \emph{static} because the single assignment may still execute many
  times, in a loop.  A program is converted by giving every definition
  of a variable a new version, \texttt{a1}, \texttt{a2}, \ldots, and
  renaming every use to the version that reaches it.
\end{definition}

\margin{LLVM IR is in SSA form for all its virtual registers (LLVM
Language Reference \cite{llvm2026langref}); the NVPTX back end of
Lesson~10 receives the program in that form.}%
In SSA form a use names its definition, so a def-use chain is a direct
link and does not need a separate analysis.  Constant propagation, dead
code elimination and most other optimisations become simpler, which is
why SSA is the basis of most compiler intermediate representations.\sidenote{LLVM
front ends do not build SSA themselves.  Clang emits every local
variable as a stack slot (\texttt{alloca}) with loads and stores, and
the \texttt{mem2reg} pass promotes such slots to SSA registers,
inserting $\phi$ functions with the standard algorithm based on the
dominance frontiers introduced below to obtain pruned SSA (LLVM passes \cite{llvm2026passes}).}

\subsection{Joins and phi functions}

Renaming is straightforward in straight-line code.  At a block with
several predecessors it is not: in \cref{fig:l11-ssa} the use of
\texttt{s} in BB4 is reached by \texttt{s1} along BB2 and by
\texttt{s2} along BB3, and no single name is correct for both paths.

\begin{definition}[Join point and phi function]
  \label{def:l11-phi}
  A \term{join point}[合流点] is a block with more than one predecessor.
  A \term{phi function}[φ関数] placed at the start of a join point,
  $x_3 = \phi(x_1, x_2)$, has one argument per predecessor and defines
  a new version $x_3$ whose value is the argument of the edge along
  which control arrived.  After the $\phi$, every later use refers to
  $x_3$.
\end{definition}

\begin{figure}[tb]
  \centering
  % Conversion to SSA form (Lesson 11).  Left: a diamond CFG in which s is
% redefined on the else path and t on the then path.  Right: the same
% CFG in SSA form, every definition a new version and two phi functions
% at the join block BB4 (arguments in the order of the predecessors
% BB2, BB3); u needs none.
% Original figure in the layout of the outline figures m11l2_ssa and
% m11l2_ssa_eg (phi drawn correctly, not as the empty-set sign).
% Usage: % Conversion to SSA form (Lesson 11).  Left: a diamond CFG in which s is
% redefined on the else path and t on the then path.  Right: the same
% CFG in SSA form, every definition a new version and two phi functions
% at the join block BB4 (arguments in the order of the predecessors
% BB2, BB3); u needs none.
% Original figure in the layout of the outline figures m11l2_ssa and
% m11l2_ssa_eg (phi drawn correctly, not as the empty-set sign).
% Usage: % Conversion to SSA form (Lesson 11).  Left: a diamond CFG in which s is
% redefined on the else path and t on the then path.  Right: the same
% CFG in SSA form, every definition a new version and two phi functions
% at the join block BB4 (arguments in the order of the predecessors
% BB2, BB3); u needs none.
% Original figure in the layout of the outline figures m11l2_ssa and
% m11l2_ssa_eg (phi drawn correctly, not as the empty-set sign).
% Usage: \input{figures/l11-ssa}
% Width: body (120 mm); the drawing is about 116 mm wide.
\begin{tikzpicture}[x=1mm, y=1mm, font=\scriptsize\sffamily,
    >={Stealth[length=1.6mm]},
    bb/.style={draw, fill=ln-box, align=left, inner sep=2pt,
               minimum width=24mm, rounded corners=0.8mm}]
  % before
  \node[font=\scriptsize\sffamily\bfseries] at (27,50) {\iflangja{変換前}{before}};
  \node[bb] (a1) at (27,40)
    {\textbf{BB1}\\\texttt{s = 1}\\\texttt{t = x + 2}\\\texttt{if (x > 0)}};
  \node[bb] (a2) at (13,24) {\textbf{BB2}\\\texttt{t = t * 3}};
  \node[bb] (a3) at (41,24) {\textbf{BB3}\\\texttt{s = 5}};
  \node[bb] (a4) at (27,7)  {\textbf{BB4}\\\texttt{u = s + t}};
  % after
  \node[font=\scriptsize\sffamily\bfseries] at (89,50) {\iflangja{SSA 形式}{SSA form}};
  \node[bb] (b1) at (89,40)
    {\textbf{BB1}\\\texttt{s1 = 1}\\\texttt{t1 = x + 2}\\\texttt{if (x > 0)}};
  \node[bb] (b2) at (75,24) {\textbf{BB2}\\\texttt{t2 = t1 * 3}};
  \node[bb] (b3) at (103,24) {\textbf{BB3}\\\texttt{s2 = 5}};
  \node[bb, fill=ln-tip!15] (b4) at (89,5)
    {\textbf{BB4}\\\texttt{s3 = $\phi$(s1, s2)}\\\texttt{t3 = $\phi$(t2, t1)}\\\texttt{u1 = s3 + t3}};
  \foreach \p in {a, b} {
    \draw[->] (\p1) -- (\p2);
    \draw[->] (\p1) -- (\p3);
    \draw[->] (\p2) -- (\p4);
    \draw[->] (\p3) -- (\p4);
  }
  \draw[->, ln-accent, very thick] (56,24) -- (61,24);
\end{tikzpicture}

% Width: body (120 mm); the drawing is about 116 mm wide.
\begin{tikzpicture}[x=1mm, y=1mm, font=\scriptsize\sffamily,
    >={Stealth[length=1.6mm]},
    bb/.style={draw, fill=ln-box, align=left, inner sep=2pt,
               minimum width=24mm, rounded corners=0.8mm}]
  % before
  \node[font=\scriptsize\sffamily\bfseries] at (27,50) {\iflangja{変換前}{before}};
  \node[bb] (a1) at (27,40)
    {\textbf{BB1}\\\texttt{s = 1}\\\texttt{t = x + 2}\\\texttt{if (x > 0)}};
  \node[bb] (a2) at (13,24) {\textbf{BB2}\\\texttt{t = t * 3}};
  \node[bb] (a3) at (41,24) {\textbf{BB3}\\\texttt{s = 5}};
  \node[bb] (a4) at (27,7)  {\textbf{BB4}\\\texttt{u = s + t}};
  % after
  \node[font=\scriptsize\sffamily\bfseries] at (89,50) {\iflangja{SSA 形式}{SSA form}};
  \node[bb] (b1) at (89,40)
    {\textbf{BB1}\\\texttt{s1 = 1}\\\texttt{t1 = x + 2}\\\texttt{if (x > 0)}};
  \node[bb] (b2) at (75,24) {\textbf{BB2}\\\texttt{t2 = t1 * 3}};
  \node[bb] (b3) at (103,24) {\textbf{BB3}\\\texttt{s2 = 5}};
  \node[bb, fill=ln-tip!15] (b4) at (89,5)
    {\textbf{BB4}\\\texttt{s3 = $\phi$(s1, s2)}\\\texttt{t3 = $\phi$(t2, t1)}\\\texttt{u1 = s3 + t3}};
  \foreach \p in {a, b} {
    \draw[->] (\p1) -- (\p2);
    \draw[->] (\p1) -- (\p3);
    \draw[->] (\p2) -- (\p4);
    \draw[->] (\p3) -- (\p4);
  }
  \draw[->, ln-accent, very thick] (56,24) -- (61,24);
\end{tikzpicture}

% Width: body (120 mm); the drawing is about 116 mm wide.
\begin{tikzpicture}[x=1mm, y=1mm, font=\scriptsize\sffamily,
    >={Stealth[length=1.6mm]},
    bb/.style={draw, fill=ln-box, align=left, inner sep=2pt,
               minimum width=24mm, rounded corners=0.8mm}]
  % before
  \node[font=\scriptsize\sffamily\bfseries] at (27,50) {\iflangja{変換前}{before}};
  \node[bb] (a1) at (27,40)
    {\textbf{BB1}\\\texttt{s = 1}\\\texttt{t = x + 2}\\\texttt{if (x > 0)}};
  \node[bb] (a2) at (13,24) {\textbf{BB2}\\\texttt{t = t * 3}};
  \node[bb] (a3) at (41,24) {\textbf{BB3}\\\texttt{s = 5}};
  \node[bb] (a4) at (27,7)  {\textbf{BB4}\\\texttt{u = s + t}};
  % after
  \node[font=\scriptsize\sffamily\bfseries] at (89,50) {\iflangja{SSA 形式}{SSA form}};
  \node[bb] (b1) at (89,40)
    {\textbf{BB1}\\\texttt{s1 = 1}\\\texttt{t1 = x + 2}\\\texttt{if (x > 0)}};
  \node[bb] (b2) at (75,24) {\textbf{BB2}\\\texttt{t2 = t1 * 3}};
  \node[bb] (b3) at (103,24) {\textbf{BB3}\\\texttt{s2 = 5}};
  \node[bb, fill=ln-tip!15] (b4) at (89,5)
    {\textbf{BB4}\\\texttt{s3 = $\phi$(s1, s2)}\\\texttt{t3 = $\phi$(t2, t1)}\\\texttt{u1 = s3 + t3}};
  \foreach \p in {a, b} {
    \draw[->] (\p1) -- (\p2);
    \draw[->] (\p1) -- (\p3);
    \draw[->] (\p2) -- (\p4);
    \draw[->] (\p3) -- (\p4);
  }
  \draw[->, ln-accent, very thick] (56,24) -- (61,24);
\end{tikzpicture}

  \caption{Conversion to SSA form.  Each definition gets a new version;
    at the join BB4, $\phi$ functions merge the versions of \texttt{s}
    and \texttt{t} that arrive from BB2 (first argument) and BB3
    (second argument).}
  \label{fig:l11-ssa}
\end{figure}

In \cref{fig:l11-ssa}, \texttt{s3 = $\phi$(s1, s2)} is \texttt{s1} when
control comes from BB2 and \texttt{s2} when it comes from BB3;
\texttt{t3} merges \texttt{t2} (BB2) and \texttt{t1} (BB3).  The use in
\texttt{u1 = s3 + t3} then has exactly one reaching definition per
name.  A $\phi$ is not an instruction of any machine.  When the program
leaves SSA form, it becomes a move on each incoming edge, or, when the
two paths are short, a select or a \term{conditional move}[条件付き移動]
that picks one of two registers by a predicate.\margin{The GPU
counterpart of the conditional move is the \texttt{selp} instruction
of PTX, which picks one of two source operands by a predicate operand
(PTX ISA \cite{nvidia2026ptx}); unlike the guarded \texttt{@p}
execution of Lesson~6 it is not itself predicated.}

LLVM writes the $\phi$ with the predecessor label beside each value.
In SSA form as \texttt{mem2reg} leaves it, before any further
optimisation, the function \texttt{min(a, b)} keeps its branch and
merges the two results in the block \texttt{end}:
\begin{lstlisting}[numbers=none]
define i32 @min(i32 %a, i32 %b) {
entry:
  %lt = icmp slt i32 %a, %b
  br i1 %lt, label %ta, label %tb
ta:
  br label %end
tb:
  br label %end
end:
  %r = phi i32 [ %a, %ta ], [ %b, %tb ]
  ret i32 %r
}
\end{lstlisting}
\margin{SimplifyCFG replaces this diamond by one
\texttt{select i1 \%lt, i32 \%a, i32 \%b}, the conditional move of the
IR (LLVM Language Reference \cite{llvm2026langref}); InstCombine then
canonicalises this min pattern to \texttt{@llvm.smin.i32}.}%
When the target is lowered, \texttt{\%r} becomes either a register
written on both incoming edges or a single conditional move.

\subsection{Where to place phi functions}

The simplest placement inserts a $\phi$ for every variable at every
join point.  It is correct but wasteful: in \cref{fig:l11-ssa} the
variable \texttt{u} is defined only in BB4 and \texttt{x} is never
redefined, so a $\phi$ for either would merge one value with itself.
A $\phi$ is needed only where different definitions actually meet.

\begin{definition}[Path-convergence criterion]
  \label{def:l11-path-convergence}
  The \term{path-convergence criterion}[経路収束条件] requires a $\phi$
  for variable $a$ at block $z$ when all of the following hold:
  \begin{enumerate}
    \item a block $x$ contains a definition of $a$;
    \item a block $y \neq x$ contains a definition of $a$;
    \item there is a nonempty path $P_{xz}$ from $x$ to $z$;
    \item there is a nonempty path $P_{yz}$ from $y$ to $z$;
    \item $P_{xz}$ and $P_{yz}$ have no node in common other than $z$;
    \item $z$ does not appear within both $P_{xz}$ and $P_{yz}$ before
      their ends, though it may appear in one of them.
  \end{enumerate}
  The start node counts as containing a definition of every variable.
  Because each inserted $\phi$ is itself a new definition of $a$, the
  criterion is applied again until no further $\phi$ is required.
\end{definition}

\caution{Lecture vs paper: production compilers do not search paths.
Cytron et al.\ place $\phi$ functions at the iterated dominance
frontier of the defining blocks, in time that is linear in practice
(quadratic in the worst case) \cite{cytron1991ssa}.  \emph{Pruned}
SSA, which \texttt{mem2reg} builds, also drops $\phi$s whose result
is never used (LLVM passes \cite{llvm2026passes}).}%
Condition~5 is the one that saves $\phi$s.  Two definitions whose paths
already met at an earlier join were merged there; the $\phi$ at that
earlier join is the single definition that flows on, and a second
$\phi$ further down would be redundant.  Condition~6 allows loops: a
path may pass through $z$ once before arriving there for the second
time along a back edge.  The start-node rule matters for a variable
assigned on one branch only: if BB3 of \cref{fig:l11-ssa} also set
\texttt{w = 7}, with no earlier assignment to \texttt{w}, the path
from the start node through BB2 and the path from BB3 would first meet
at BB4, so BB4 would receive \texttt{w2 = $\phi$(w0, w1)}, whose first
argument is the initial version \texttt{w0} of the start node.

\begin{figure}[tb]
  \centering
  % Where phi functions go in a loop (Lesson 11): the SSA form of a
% counting loop.  B2 is a join of the entry path and the back edge, so
% i and s receive phi functions there; n has one definition and needs
% none; B4 has a single predecessor, so no phi is placed there although
% two definitions of s reach it.
% Original figure (applies the criterion of the outline figure
% m11l2_path_converge to the book's own CFG).
% Usage: % Where phi functions go in a loop (Lesson 11): the SSA form of a
% counting loop.  B2 is a join of the entry path and the back edge, so
% i and s receive phi functions there; n has one definition and needs
% none; B4 has a single predecessor, so no phi is placed there although
% two definitions of s reach it.
% Original figure (applies the criterion of the outline figure
% m11l2_path_converge to the book's own CFG).
% Usage: % Where phi functions go in a loop (Lesson 11): the SSA form of a
% counting loop.  B2 is a join of the entry path and the back edge, so
% i and s receive phi functions there; n has one definition and needs
% none; B4 has a single predecessor, so no phi is placed there although
% two definitions of s reach it.
% Original figure (applies the criterion of the outline figure
% m11l2_path_converge to the book's own CFG).
% Usage: \input{figures/l11-path-convergence}
% Width: body (120 mm); the drawing is about 112 mm wide.
\begin{tikzpicture}[x=1mm, y=1mm, font=\scriptsize\sffamily,
    >={Stealth[length=1.6mm]},
    bb/.style={draw, fill=ln-box, align=left, inner sep=2pt,
               minimum width=30mm, rounded corners=0.8mm}]
  \node[bb] (b1) at (22,46) {\textbf{B1}\\\texttt{i1 = 0}\\\texttt{s1 = 0}};
  \node[bb, fill=ln-tip!15] (b2) at (22,27)
    {\textbf{B2}\\\texttt{i2 = $\phi$(i1, i3)}\\\texttt{s2 = $\phi$(s1, s3)}\\\texttt{if (i2 < n)}};
  \node[bb] (b3) at (66,27) {\textbf{B3}\\\texttt{s3 = s2 + i2}\\\texttt{i3 = i2 + 1}};
  \node[bb] (b4) at (22,6) {\textbf{B4}\\\texttt{return s2}};
  \draw[->] (b1) -- (b2);
  \draw[->] (b2.east) -- node[above] {\texttt{i2 < n}} (b3.west);
  \draw[->, ln-caution, thick] (b3.south) -- ++(0,-6) -| 
    node[pos=0.25, below] {\iflangja{後退辺}{back edge}} ([xshift=4mm]b2.south);
  \draw[->] ([xshift=-6mm]b2.south) -- node[left] {\texttt{i2 >= n}} ([xshift=-6mm]b4.north);
  \node[anchor=west, align=left, text=ln-muted, text width=34mm] at (84,40)
    {\iflangja{B2: 2経路（B1 から，B3 から）が B2 でのみ交わる $\Rightarrow$ $\phi$}{B2: the paths from B1 and from B3 meet only at B2 $\Rightarrow$ $\phi$}};
  \node[anchor=west, align=left, text=ln-muted, text width=34mm] at (84,8)
    {\iflangja{B4: 経路はすべて B2 を通る $\Rightarrow$ $\phi$ 不要}{B4: every path passes through B2 $\Rightarrow$ no $\phi$}};
\end{tikzpicture}

% Width: body (120 mm); the drawing is about 112 mm wide.
\begin{tikzpicture}[x=1mm, y=1mm, font=\scriptsize\sffamily,
    >={Stealth[length=1.6mm]},
    bb/.style={draw, fill=ln-box, align=left, inner sep=2pt,
               minimum width=30mm, rounded corners=0.8mm}]
  \node[bb] (b1) at (22,46) {\textbf{B1}\\\texttt{i1 = 0}\\\texttt{s1 = 0}};
  \node[bb, fill=ln-tip!15] (b2) at (22,27)
    {\textbf{B2}\\\texttt{i2 = $\phi$(i1, i3)}\\\texttt{s2 = $\phi$(s1, s3)}\\\texttt{if (i2 < n)}};
  \node[bb] (b3) at (66,27) {\textbf{B3}\\\texttt{s3 = s2 + i2}\\\texttt{i3 = i2 + 1}};
  \node[bb] (b4) at (22,6) {\textbf{B4}\\\texttt{return s2}};
  \draw[->] (b1) -- (b2);
  \draw[->] (b2.east) -- node[above] {\texttt{i2 < n}} (b3.west);
  \draw[->, ln-caution, thick] (b3.south) -- ++(0,-6) -| 
    node[pos=0.25, below] {\iflangja{後退辺}{back edge}} ([xshift=4mm]b2.south);
  \draw[->] ([xshift=-6mm]b2.south) -- node[left] {\texttt{i2 >= n}} ([xshift=-6mm]b4.north);
  \node[anchor=west, align=left, text=ln-muted, text width=34mm] at (84,40)
    {\iflangja{B2: 2経路（B1 から，B3 から）が B2 でのみ交わる $\Rightarrow$ $\phi$}{B2: the paths from B1 and from B3 meet only at B2 $\Rightarrow$ $\phi$}};
  \node[anchor=west, align=left, text=ln-muted, text width=34mm] at (84,8)
    {\iflangja{B4: 経路はすべて B2 を通る $\Rightarrow$ $\phi$ 不要}{B4: every path passes through B2 $\Rightarrow$ no $\phi$}};
\end{tikzpicture}

% Width: body (120 mm); the drawing is about 112 mm wide.
\begin{tikzpicture}[x=1mm, y=1mm, font=\scriptsize\sffamily,
    >={Stealth[length=1.6mm]},
    bb/.style={draw, fill=ln-box, align=left, inner sep=2pt,
               minimum width=30mm, rounded corners=0.8mm}]
  \node[bb] (b1) at (22,46) {\textbf{B1}\\\texttt{i1 = 0}\\\texttt{s1 = 0}};
  \node[bb, fill=ln-tip!15] (b2) at (22,27)
    {\textbf{B2}\\\texttt{i2 = $\phi$(i1, i3)}\\\texttt{s2 = $\phi$(s1, s3)}\\\texttt{if (i2 < n)}};
  \node[bb] (b3) at (66,27) {\textbf{B3}\\\texttt{s3 = s2 + i2}\\\texttt{i3 = i2 + 1}};
  \node[bb] (b4) at (22,6) {\textbf{B4}\\\texttt{return s2}};
  \draw[->] (b1) -- (b2);
  \draw[->] (b2.east) -- node[above] {\texttt{i2 < n}} (b3.west);
  \draw[->, ln-caution, thick] (b3.south) -- ++(0,-6) -| 
    node[pos=0.25, below] {\iflangja{後退辺}{back edge}} ([xshift=4mm]b2.south);
  \draw[->] ([xshift=-6mm]b2.south) -- node[left] {\texttt{i2 >= n}} ([xshift=-6mm]b4.north);
  \node[anchor=west, align=left, text=ln-muted, text width=34mm] at (84,40)
    {\iflangja{B2: 2経路（B1 から，B3 から）が B2 でのみ交わる $\Rightarrow$ $\phi$}{B2: the paths from B1 and from B3 meet only at B2 $\Rightarrow$ $\phi$}};
  \node[anchor=west, align=left, text=ln-muted, text width=34mm] at (84,8)
    {\iflangja{B4: 経路はすべて B2 を通る $\Rightarrow$ $\phi$ 不要}{B4: every path passes through B2 $\Rightarrow$ no $\phi$}};
\end{tikzpicture}

  \caption{SSA form of a counting loop.  The loop header B2 joins the
    entry path and the back edge, so \texttt{i} and \texttt{s} get
    $\phi$ functions there and nowhere else.}
  \label{fig:l11-path-convergence}
\end{figure}

\begin{example}
  \label{ex:l11-path-convergence}
  \Cref{fig:l11-path-convergence} sums the integers below \texttt{n}.
  Before renaming, \texttt{i} and \texttt{s} are defined in B1 and in B3.
  The paths B1$\to$B2 and B3$\to$B2 share only B2, so both conditions
  5 and 6 hold at B2 and each variable gets a $\phi$ there; its second
  argument comes from the back edge.  At B4 every path from B1 and
  every path from B3 runs through B2, so condition~5 fails and B4 uses
  \texttt{s2} directly, although two definitions of \texttt{s} reach
  it.  \texttt{n} is defined only at the start, so condition~2 fails
  and it needs no $\phi$ anywhere.  The iteration adds nothing either:
  the new definitions in B2 meet the others only at B2 itself, which
  already has its $\phi$s.
\end{example}

The faster algorithm rests on dominance, the forward counterpart of the
post-dominance of Lesson~6.  A block $d$ is a
\term{dominator}[支配節点] of a block $n$ if every path from the entry
to $n$ passes through $d$.  The \term{dominance frontier}[支配辺境] of
$d$ is the set of blocks $z$ such that $d$ dominates a predecessor of
$z$ but does not strictly dominate $z$ itself: the first blocks at
which paths from $d$ meet paths that bypass $d$
\cite{cytron1991ssa}.\margin{In \cref{fig:l11-path-convergence} B2
dominates B3 and B4; the dominance frontier of B3 is $\{\mathrm{B2}\}$,
exactly where its definitions need $\phi$s.  Dominators: Aho et al.\
§9.6 \cite{aho2006compilers}.}  A definition in $d$ needs $\phi$
functions exactly at the blocks of the frontier of $d$, then at the
frontiers of those blocks, and so on, which is the result of the
path-convergence criterion without enumerating paths.

\begin{keypoint}
  SSA gives every assignment its own name, so each use has one
  definition.  At a join point, $\phi$ functions merge the versions
  arriving along different edges; they are needed only where distinct
  definitions first meet, which the path-convergence criterion (or
  dominance frontiers) determines.  A $\phi$ is lowered to moves or a
  conditional move.
\end{keypoint}

\section{Examples of compiler optimizations}
\label{sec:l11-optimizations}

The analyses above exist to justify transformations.\source{M11,
compiler optimizations; Aho et al.\ §9.1 The Principal Sources of
Optimization \cite{aho2006compilers}; CPG §7.39 \#pragma unroll, §7.1.5
\cite{nvidia2024cudaguide}.}  Each optimisation below must preserve the
meaning of the program, and each needs some fact about it that a
data-flow analysis or the SSA form provides.  The examples are given in
CUDA C++ or PTX; \cref{tab:l11-optimizations} summarises them.

\subsection{Loop unrolling}

\begin{definition}[Loop unrolling]
  \label{def:l11-unroll}
  The transformation called \term{loop unrolling}[ループ展開] replicates
  the body of a loop $u$
  times, the \emph{unroll factor}, so that one iteration of the new
  loop does the work of $u$ old iterations.  When the trip count is
  unknown or not a multiple of $u$, an
  \term{epilogue loop}[エピローグループ] after the unrolled loop
  executes the remaining iterations one at a time.
\end{definition}

\begin{lstlisting}[language=C++, numbers=none]
// before
for (int i = 0; i < n; i++)
  y[i] = c * x[i];

// unrolled by 4, with an epilogue loop
int i = 0;
for (; i + 3 < n; i += 4) {
  y[i]     = c * x[i];
  y[i + 1] = c * x[i + 1];
  y[i + 2] = c * x[i + 2];
  y[i + 3] = c * x[i + 3];
}
for (; i < n; i++)          // 0 to 3 iterations
  y[i] = c * x[i];
\end{lstlisting}

\margin{\enquote{By default, the compiler unrolls small loops with a
known trip count}; \texttt{\#pragma unroll} with an optional factor
controls any loop, and \texttt{\#pragma unroll 1} forbids unrolling
(CPG §7.39 \cite{nvidia2024cudaguide}).}%
Unrolling removes three of every four compare-and-branch sequences and
increment instructions.  More importantly, it puts four independent
multiplications and memory accesses into one basic block, where the
scheduler can interleave them to hide latency and a vectorising
compiler can combine them.  The cost is code size and, on a GPU,
registers: the four loaded values are live at the same time.\caution{An
unroll factor that pushes a kernel over its register budget can cost
more through lost occupancy or spills than it saves in branches.}

\subsection{Function inlining}

\begin{definition}[Function inlining]
  \label{def:l11-inlining}
  The transformation called \term{function inlining}[関数のインライン展開]
  replaces a call by a copy
  of the body of the called function, with the arguments substituted
  for the parameters.  Deciding it needs facts about both caller and
  callee, so it is an \term{inter-procedural analysis}[手続き間解析]
  decision, as opposed to the intra-procedural analyses of the previous
  sections.
\end{definition}

\begin{lstlisting}[language=C++, numbers=none]
__device__ float sq(float v) { return v * v; }
__device__ float norm2(float a, float b) {
  return sq(a) + sq(b);
}
// after inlining both levels at a call norm2(x, y):
float r = x * x + y * y;
\end{lstlisting}

Inlining removes the call and return, the passing of arguments and the
saving of registers on the stack.  It also exposes the body to the
optimisations of the caller: constants of the call site propagate into
it, and code that is dead for these arguments disappears.  The cost is
code size, since every call site receives its own copy.\margin{The
CUDA compiler inlines a \texttt{\_\_device\_\_} function \enquote{when
deemed appropriate}; \texttt{\_\_noinline\_\_} asks it not to and
\texttt{\_\_forceinline\_\_} forces it (CPG §7.1.5
\cite{nvidia2024cudaguide}).}

\subsection{Dead code elimination}

\begin{definition}[Dead code elimination]
  \label{def:l11-dce}
  The optimisation called \term{dead code elimination}[不要コード除去]
  removes instructions whose
  results never influence the output of the program: an assignment to a
  variable that is dead immediately after it, and code that can never
  execute.
\end{definition}

\begin{lstlisting}[language=PTX, numbers=none]
ld.param.f32  %f1, [x];
mul.f32       %f2, %f1, %f1;   // dead: %f2 is
add.f32       %f2, %f1, 1.0;   //   rewritten here
st.global.f32 [%rd1], %f2;
\end{lstlisting}

The product written by \texttt{mul.f32} is overwritten before any
instruction reads it: \texttt{\%f2} is dead after the \texttt{mul}, in
the sense of \cref{def:l11-live}, so the instruction is removed.  The
argument is a liveness argument; in SSA form it becomes simpler still,
because a definition with no uses is dead by inspection.  Removing an
instruction can make the definitions it read dead in turn, so the
elimination is repeated.\caution{An instruction with a side effect (a
store, an atomic, a barrier) is never dead merely because its result is
unused; only pure computations are removed.}

\subsection{Constant propagation and folding}

\begin{definition}[Constant propagation and folding]
  \label{def:l11-constprop}
  The optimisation called \term{constant propagation}[定数伝播] replaces
  a use of a variable by a constant when every definition reaching the
  use assigns that constant.  The optimisation called
  \term{constant folding}[定数畳み込み] evaluates at compile time an
  operation all of whose operands are constants.
\end{definition}

\begin{lstlisting}[language=PTX, numbers=none]
// before
mov.f32  %f1, 2.0;
mov.f32  %f2, 4.0;
add.f32  %f3, %f1, %f2;
mul.f32  %f4, %f3, %f5;
// after propagation, folding, propagation
mul.f32  %f4, %f5, 6.0;
\end{lstlisting}

Propagating the two constants into the \texttt{add} makes it
$2.0 + 4.0$, which folding turns into the constant $6.0$; propagating
that constant into the \texttt{mul} leaves the two \texttt{mov}
instructions and the \texttt{add} without uses, and dead code
elimination removes them.  One optimisation thus enables the
next.\margin{Constant propagation is a client of reaching definitions
(Lesson~10): a use is constant if all definitions reaching it assign
the same constant.  Copy propagation, also from Lesson~10, works the
same way for copies.}

\subsection{Strength reduction}

\begin{definition}[Strength reduction]
  \label{def:l11-strength}
  The optimisation called \term{strength reduction}[強さの低減] replaces
  an operation by an
  equivalent one that costs less: a division by a multiplication, a
  multiplication by a shift or an addition, a call to a power function
  by a multiplication.
\end{definition}

The two classic examples are a division by two, which becomes a right
shift, and a square, which becomes a multiplication of the value by
itself.  In PTX they are given as the following pairs, the square in
the form of an \texttt{ex2.approx.f32}:
\begin{lstlisting}[language=PTX, numbers=none]
div.rn.f32     %f2, %f1, 2.0;   // becomes
mul.f32        %f2, %f1, 0.5;
ex2.approx.f32 %f3, %f1;        // becomes
mul.f32        %f3, %f1, %f1;
\end{lstlisting}
\caution{Lecture vs ISA: both pairs are mislabelled.  The first is a
floating-point division turned into a multiplication, not a shift; a
right shift replaces division only for unsigned or non-negative
integers, since for a negative signed value it rounds toward $-\infty$
while division rounds toward zero.  And \texttt{ex2} computes $2^x$,
not $x^2$ (PTX ISA \cite{nvidia2026ptx}), so the second pair is not
an equivalence; the correct square example is \texttt{pow(a, 2)}
$\to$ \texttt{a * a}.}%
A floating-point division by two equals a multiplication by $0.5$
exactly, because both scale by a power of two.  The shift belongs to
integers: an unsigned division by $2^k$, such as \texttt{div.u32 \%r2,
\%r1, 8}, becomes the right shift \texttt{shr.u32 \%r2, \%r1, 3}, while
a signed one needs a correction for negative values.  A square is
reduced from a call: \texttt{pow(a, 2)} becomes \texttt{a * a}.
Division by a constant that is not a
power of two, $x/3$, cannot be replaced by $x\cdot(1/3)$ in floating
point without changing some results, so the compiler does it only when
the program allows approximate division, as \texttt{nvcc
--prec-div=false} (implied by \texttt{-use\_fast\_math}) does for
single precision.  A divisor that is a power of two only at run time is
beyond the compiler; Lesson~5 shows the rewrite by hand.

\subsection{Loop-invariant code motion and induction variables}

\begin{definition}[Code motion]
  \label{def:l11-code-motion}
  The transformation called \term{code motion}[コード移動] moves a
  computation to another place in
  the program where it executes less often with the same result.  Its
  most common form, \term{loop-invariant code motion}[ループ不変式の移動],
  hoists a computation whose operands do not change inside a loop out of
  the loop, so that it executes once instead of once per iteration.
\end{definition}

\begin{lstlisting}[language=C++, numbers=none]
// before: scale is recomputed in every iteration
for (int i = 0; i < n; i++) {
  float scale = alpha / beta;
  y[i] += scale * x[i];
}
// after: computed once, before the loop
float scale = alpha / beta;
for (int i = 0; i < n; i++)
  y[i] += scale * x[i];
\end{lstlisting}

The number of dynamic instructions falls by one division per
iteration.\margin{An instruction is loop invariant when each of its
operands is a constant, is defined outside the loop, or is defined by
another loop-invariant instruction; the test is applied repeatedly
(Aho et al.\ §9.1 \cite{aho2006compilers}).}  To be allowed, the compiler must know that neither
\texttt{alpha} nor \texttt{beta} is changed in the loop, a question for
reaching definitions, and that executing the division when the loop
runs zero times is harmless.

\begin{definition}[Induction variable]
  \label{def:l11-induction}
  An \term{induction variable}[誘導変数] is a variable that changes by
  the same amount in every iteration of a loop, such as the counter
  \texttt{i} that controls it.  In PTX its update is a single
  instruction such as \texttt{add.s32 \%r4, \%r4, 1;}.
\end{definition}

Recognising induction variables enables further optimisations.  The
address of \texttt{x[i]} is $\mathit{base} + 4i$, itself an induction
variable that grows by $4$ per iteration, so strength reduction can
replace the multiplication by an addition of $4$ to a pointer; the
trip count of the loop, needed for unrolling, follows from the initial
value, step and bound of the counter.

\begin{table}[tb]
  \centering\small
  \begin{tabular}{@{}>{\raggedright\arraybackslash}p{30mm}
                    >{\raggedright\arraybackslash}p{42mm}
                    >{\raggedright\arraybackslash}p{38mm}@{}}
    \toprule
    Optimisation & What it saves & What it needs to know \\
    \midrule
    Loop unrolling & branches, counter updates; exposes parallelism &
      trip count, register budget \\
    Function inlining & call overhead; enables other optimisations &
      callee body (inter-procedural) \\
    Dead code elimination & instructions with unused results & liveness
      or SSA uses \\
    Constant propagation and folding & run-time computation of constants
      & reaching definitions \\
    Strength reduction & expensive operations & exact equivalence of the
      cheaper form \\
    Loop-invariant code motion & repeated work in loops & reaching
      definitions; loop structure \\
    \bottomrule
  \end{tabular}
  \caption{The optimisations of this section and the analyses behind
    them.}
  \label{tab:l11-optimizations}
\end{table}

\section{Divergence analysis}
\label{sec:l11-divergence}

Lesson~6 showed how the hardware executes a divergent branch: the SIMT
stack runs the paths one after the other under partial masks and joins
them at the reconvergence point, normally the immediate post-dominator
of the branch.\source{M11, divergence analysis; CPG §5.4.2 Control
Flow Instructions, §4.1 SIMT Architecture \cite{nvidia2024cudaguide};
PMPP ch.~10 \cite{hwu2022pmpp}.}  The compiler supplies the
reconvergence information, and it can omit the stack operations, or
emit simpler code, for branches it proves warp-uniform.  It therefore
needs to know statically which branches may diverge.\margin{PTX can
record the result: \texttt{bra.uni} marks a branch that is guaranteed
to be non-divergent, all active threads of the warp having the same
predicate and target (PTX ISA, \texttt{bra} \cite{nvidia2026ptx}).}

\begin{definition}[Divergence analysis]
  \label{def:l11-divergence-analysis}
  The static analysis called
  \term{divergence analysis}[ダイバージェンス解析] is one
  that determines, for each branch of a kernel, whether its condition
  may evaluate differently in different threads of a warp.  Its output
  is the set of branches that may diverge and need reconvergence
  information; every other branch is uniform.
\end{definition}

The question predates GPUs.  A compiler that turns a loop into vector
instructions, \term{vectorization}[ベクトル化], must check that all
iterations it packs together execute the same instructions in lockstep,
including inside branches of the loop body; GPU compilers inherited
these tests.\margin{The two settings are the same problem: the
iterations of a vectorised loop play the role of the threads of a warp,
and a branch on the loop index plays the role of a branch on the
thread index.}  For a simple if-else the test is easy; with complex
control flow and synchronisation it becomes hard.

\subsection{Thread-dependent and uniform values}

\begin{definition}[Thread-dependent and uniform values]
  \label{def:l11-thread-dependent}
  An expression is a \term{thread-dependent value}[スレッド依存値] at a
  program point $p$ if it may evaluate to different values in different
  threads at $p$.  It is a \term{uniform value}[一様値] at $p$ if it has
  the same value in all threads of a warp (or, more strongly, of the
  kernel).  A branch is divergent only if its condition is thread
  dependent.
\end{definition}

The source of thread dependence is the thread index: \texttt{threadIdx}
differs in every lane.  Kernel parameters, \texttt{blockIdx},
\texttt{blockDim} and constants are uniform within a block, and so are
values computed only from them.  Dependence then spreads along the
def-use chains: a value computed from a thread-dependent operand is
thread dependent.  The kernel below has four branches.

\begin{lstlisting}[language=C++, numbers=none]
__global__ void k(const int *mem, int *out, int n) {
  int tid  = threadIdx.x;
  int q    = tid / 2;
  int base = blockIdx.x * n;
  int x    = mem[base];        // same address in all
  int y    = mem[base + tid];  // address per thread
  int acc  = 0;
  if (q > 3)  acc += 1;        // divergent
  if (y > 0)  acc += 2;        // divergent
  if (x > 0)  acc += 4;        // uniform
  for (int j = 0; j < n; j++)  // uniform
    acc += j;
  out[base + tid] = acc;
}
\end{lstlisting}

\Cref{fig:l11-divergence} draws the def-use edges.  The condition
\texttt{q > 3} depends, through \texttt{tid} and \texttt{q}, on
\texttt{threadIdx.x}: indirect dependence is still dependence.  The two loads
differ in their addresses.  A value loaded from memory is thread
dependent if its address is, because different threads then read
different words; \texttt{mem[base]} is read at one address by all
threads, so \texttt{x} is uniform and so is the branch on it.\margin{Uniform
within a warp is all that divergence needs: \texttt{blockIdx.x} and
\texttt{base} differ between blocks, but never between the threads of
one warp, since a warp never spans two blocks.}  The loop condition \texttt{j < n} depends
only on the parameter \texttt{n} and the uniform counter \texttt{j}:
all threads run the same number of iterations.

\begin{figure}[tb]
  \centering
  % Divergence analysis as dependence propagation (Lesson 11): the
% def-use edges of the kernel in the text, from the sources threadIdx.x
% (thread dependent) and blockIdx.x, n (uniform) to the four branch
% conditions.  Red nodes are thread dependent, blue nodes uniform; the
% dashed edge is a dependence through a load address.
% Original figure (the idea of the outline figures m11l4_td and
% m11l4_check, on the book's own kernel).
% Usage: % Divergence analysis as dependence propagation (Lesson 11): the
% def-use edges of the kernel in the text, from the sources threadIdx.x
% (thread dependent) and blockIdx.x, n (uniform) to the four branch
% conditions.  Red nodes are thread dependent, blue nodes uniform; the
% dashed edge is a dependence through a load address.
% Original figure (the idea of the outline figures m11l4_td and
% m11l4_check, on the book's own kernel).
% Usage: % Divergence analysis as dependence propagation (Lesson 11): the
% def-use edges of the kernel in the text, from the sources threadIdx.x
% (thread dependent) and blockIdx.x, n (uniform) to the four branch
% conditions.  Red nodes are thread dependent, blue nodes uniform; the
% dashed edge is a dependence through a load address.
% Original figure (the idea of the outline figures m11l4_td and
% m11l4_check, on the book's own kernel).
% Usage: \input{figures/l11-divergence}
% Width: body (120 mm); the drawing is about 116 mm wide.
\begin{tikzpicture}[x=1mm, y=1mm, font=\scriptsize\sffamily,
    >={Stealth[length=1.6mm]},
    v/.style={draw, rounded corners=0.8mm, inner sep=2pt,
              minimum height=5.5mm, font=\scriptsize\ttfamily},
    var/.style={v, fill=ln-caution!18},
    uni/.style={v, fill=ln-accent!15},
    br/.style={v, very thick}]
  \node[font=\scriptsize\sffamily\bfseries] at (8,52) {\iflangja{発生源}{sources}};
  \node[font=\scriptsize\sffamily\bfseries] at (50,52) {\iflangja{値}{values}};
  \node[font=\scriptsize\sffamily\bfseries] at (100,52) {\iflangja{分岐条件}{branch conditions}};
  \node[var] (tidx) at (8,40) {threadIdx.x};
  \node[uni] (bidx) at (8,22) {blockIdx.x};
  \node[uni] (n)    at (8,6)  {n};
  \node[var] (tid)  at (34,40) {tid};
  \node[var] (q)    at (58,46) {q = tid / 2};
  \node[uni] (base) at (34,22) {base};
  \node[var] (y)    at (62,32) {y = mem[base + tid]};
  \node[uni] (x)    at (62,16) {x = mem[base]};
  \node[br, draw=ln-caution, fill=ln-caution!18] (c1) at (102,46) {q > 3};
  \node[br, draw=ln-caution, fill=ln-caution!18] (c2) at (102,32) {y > 0};
  \node[br, draw=ln-accent, fill=ln-accent!15] (c3) at (102,16) {x > 0};
  \node[br, draw=ln-accent, fill=ln-accent!15] (c4) at (102,4)  {j < n};
  \draw[->] (tidx) -- (tid);
  \draw[->] (tid) -- (q);
  \draw[->, dashed] (tid) -- node[pos=0.45, below left, inner sep=1pt,
    font=\scriptsize\sffamily] {\iflangja{アドレス}{address}} (y);
  \draw[->] (bidx) -- (base);
  \draw[->] (n) -- (base);
  \draw[->] (base) -- (y);
  \draw[->] (base) -- (x);
  \draw[->] (q) -- (c1);
  \draw[->] (y) -- (c2);
  \draw[->] (x) -- (c3);
  \draw[->] (n) -- (c4);
  \node[anchor=west, text=ln-caution] at (112,46) {\iflangja{発散}{div.}};
  \node[anchor=west, text=ln-caution] at (112,32) {\iflangja{発散}{div.}};
  \node[anchor=west, text=ln-accent] at (112,16) {\iflangja{一様}{unif.}};
  \node[anchor=west, text=ln-accent] at (112,4)  {\iflangja{一様}{unif.}};
\end{tikzpicture}

% Width: body (120 mm); the drawing is about 116 mm wide.
\begin{tikzpicture}[x=1mm, y=1mm, font=\scriptsize\sffamily,
    >={Stealth[length=1.6mm]},
    v/.style={draw, rounded corners=0.8mm, inner sep=2pt,
              minimum height=5.5mm, font=\scriptsize\ttfamily},
    var/.style={v, fill=ln-caution!18},
    uni/.style={v, fill=ln-accent!15},
    br/.style={v, very thick}]
  \node[font=\scriptsize\sffamily\bfseries] at (8,52) {\iflangja{発生源}{sources}};
  \node[font=\scriptsize\sffamily\bfseries] at (50,52) {\iflangja{値}{values}};
  \node[font=\scriptsize\sffamily\bfseries] at (100,52) {\iflangja{分岐条件}{branch conditions}};
  \node[var] (tidx) at (8,40) {threadIdx.x};
  \node[uni] (bidx) at (8,22) {blockIdx.x};
  \node[uni] (n)    at (8,6)  {n};
  \node[var] (tid)  at (34,40) {tid};
  \node[var] (q)    at (58,46) {q = tid / 2};
  \node[uni] (base) at (34,22) {base};
  \node[var] (y)    at (62,32) {y = mem[base + tid]};
  \node[uni] (x)    at (62,16) {x = mem[base]};
  \node[br, draw=ln-caution, fill=ln-caution!18] (c1) at (102,46) {q > 3};
  \node[br, draw=ln-caution, fill=ln-caution!18] (c2) at (102,32) {y > 0};
  \node[br, draw=ln-accent, fill=ln-accent!15] (c3) at (102,16) {x > 0};
  \node[br, draw=ln-accent, fill=ln-accent!15] (c4) at (102,4)  {j < n};
  \draw[->] (tidx) -- (tid);
  \draw[->] (tid) -- (q);
  \draw[->, dashed] (tid) -- node[pos=0.45, below left, inner sep=1pt,
    font=\scriptsize\sffamily] {\iflangja{アドレス}{address}} (y);
  \draw[->] (bidx) -- (base);
  \draw[->] (n) -- (base);
  \draw[->] (base) -- (y);
  \draw[->] (base) -- (x);
  \draw[->] (q) -- (c1);
  \draw[->] (y) -- (c2);
  \draw[->] (x) -- (c3);
  \draw[->] (n) -- (c4);
  \node[anchor=west, text=ln-caution] at (112,46) {\iflangja{発散}{div.}};
  \node[anchor=west, text=ln-caution] at (112,32) {\iflangja{発散}{div.}};
  \node[anchor=west, text=ln-accent] at (112,16) {\iflangja{一様}{unif.}};
  \node[anchor=west, text=ln-accent] at (112,4)  {\iflangja{一様}{unif.}};
\end{tikzpicture}

% Width: body (120 mm); the drawing is about 116 mm wide.
\begin{tikzpicture}[x=1mm, y=1mm, font=\scriptsize\sffamily,
    >={Stealth[length=1.6mm]},
    v/.style={draw, rounded corners=0.8mm, inner sep=2pt,
              minimum height=5.5mm, font=\scriptsize\ttfamily},
    var/.style={v, fill=ln-caution!18},
    uni/.style={v, fill=ln-accent!15},
    br/.style={v, very thick}]
  \node[font=\scriptsize\sffamily\bfseries] at (8,52) {\iflangja{発生源}{sources}};
  \node[font=\scriptsize\sffamily\bfseries] at (50,52) {\iflangja{値}{values}};
  \node[font=\scriptsize\sffamily\bfseries] at (100,52) {\iflangja{分岐条件}{branch conditions}};
  \node[var] (tidx) at (8,40) {threadIdx.x};
  \node[uni] (bidx) at (8,22) {blockIdx.x};
  \node[uni] (n)    at (8,6)  {n};
  \node[var] (tid)  at (34,40) {tid};
  \node[var] (q)    at (58,46) {q = tid / 2};
  \node[uni] (base) at (34,22) {base};
  \node[var] (y)    at (62,32) {y = mem[base + tid]};
  \node[uni] (x)    at (62,16) {x = mem[base]};
  \node[br, draw=ln-caution, fill=ln-caution!18] (c1) at (102,46) {q > 3};
  \node[br, draw=ln-caution, fill=ln-caution!18] (c2) at (102,32) {y > 0};
  \node[br, draw=ln-accent, fill=ln-accent!15] (c3) at (102,16) {x > 0};
  \node[br, draw=ln-accent, fill=ln-accent!15] (c4) at (102,4)  {j < n};
  \draw[->] (tidx) -- (tid);
  \draw[->] (tid) -- (q);
  \draw[->, dashed] (tid) -- node[pos=0.45, below left, inner sep=1pt,
    font=\scriptsize\sffamily] {\iflangja{アドレス}{address}} (y);
  \draw[->] (bidx) -- (base);
  \draw[->] (n) -- (base);
  \draw[->] (base) -- (y);
  \draw[->] (base) -- (x);
  \draw[->] (q) -- (c1);
  \draw[->] (y) -- (c2);
  \draw[->] (x) -- (c3);
  \draw[->] (n) -- (c4);
  \node[anchor=west, text=ln-caution] at (112,46) {\iflangja{発散}{div.}};
  \node[anchor=west, text=ln-caution] at (112,32) {\iflangja{発散}{div.}};
  \node[anchor=west, text=ln-accent] at (112,16) {\iflangja{一様}{unif.}};
  \node[anchor=west, text=ln-accent] at (112,4)  {\iflangja{一様}{unif.}};
\end{tikzpicture}

  \caption{The def-use edges of the kernel above, from the sources to
    the branch conditions.  Thread dependence (red) spreads from
    \texttt{threadIdx.x}; the condition of a branch decides whether it
    may diverge.}
  \label{fig:l11-divergence}
\end{figure}

Checking a branch thus becomes a reachability question on the def-use
chains: does a definition that depends on a thread index reach the
condition?  As with reaching definitions, the dependence facts are
propagated through the control-flow graph until nothing changes; loops
make the iteration necessary, because a value can become thread
dependent through its own back edge.\caution{The lecture tracks data
dependence only, which is not enough for a sound analysis.  After \texttt{b = 0; if (q >
3) b = 1;} the value of \texttt{b} at the join differs between threads
although both assignments are constants, because the branch that chose
between them diverged.  A sound analysis also marks values defined
under a divergent branch and merged at its join (a $\phi$ in SSA form)
as thread dependent: \emph{control dependence}.}

\subsection{Conservative analysis}

\begin{definition}[Conservative analysis]
  \label{def:l11-conservative}
  A static analysis is a \term{conservative analysis}[保守的解析] when
  its approximation errs only on the safe side: here, it may report a
  uniform branch as divergent but never a divergent branch as uniform.
\end{definition}

Treating a uniform branch as divergent costs only performance: an
unnecessary reconvergence point and the stack operations for it.\margin{CPG §5.4.2 gives the
standard uniform branch that a naive analysis calls divergent: a condition that depends only on
\texttt{threadIdx / warpSize} cannot diverge, since warps consist of
consecutive thread IDs (CPG §4.1 \cite{nvidia2024cudaguide}), yet it is
derived from \texttt{threadIdx}.}%
Treating a divergent branch as uniform would be wrong, so every
practical analysis is conservative.  The difficulty is precision: an
analysis that marks every value derived from \texttt{threadIdx} as
thread dependent soon marks all branches and all loops as possibly
divergent, and the analysis is useless.  Compilers therefore add rules that
prune such false positives, for example by recognising values that are
uniform within a warp although they differ between the warps of a
block (such as \texttt{threadIdx.x / warpSize}), and the larger
and more irregular the control-flow graph, the harder it is to stay
both safe and precise.\tip{Writing a branch condition in terms of
\texttt{threadIdx.x / 32} (in a one-dimensional block),
\texttt{blockIdx} or kernel parameters keeps it warp-uniform; this is the programmer's side of the analysis
(Lesson~5).}

\begin{keypoint}[Lesson summary]
  Liveness is the backward data-flow problem: IN$[B] = \mathit{use}_B
  \cup (\mathrm{OUT}[B] - \mathit{def}_B)$, OUT$[B]$ the union of the
  successors' IN sets, IN[EXIT] $= \varnothing$, solved iteratively to a
  fixed point.  Register allocation colours the interference graph of
  live ranges; when registers run out, \texttt{ptxas} spills to local
  memory, trading occupancy against memory traffic.  SSA form gives
  every definition its own name and merges versions with $\phi$
  functions at join points where distinct definitions first meet.
  Unrolling, inlining, dead code elimination, constant propagation and
  folding, strength reduction and loop-invariant code motion each rest
  on such an analysis.  Divergence analysis propagates thread
  dependence from \texttt{threadIdx} along def-use chains (and, to be
  sound, control dependence) to the branch conditions; it must be
  conservative, and its value lies in pruning the uniform branches.
\end{keypoint}

\end{document}
